\chapter*{Agradecimentos}

\addcontentsline{toc}{chapter}{Agradecimentos}

Começo por agradecer ao Professor João Pires pelo apoio e pela paciência durante todo este 
processo, e em particular pelas várias rondas de revisões detalhadas de cada capítulo desta dissertação, 
sempre feitas com rigor e com sugestões construtivas. Desde o momento em que aceitou ser meu orientador que tem sido alguém com 
quem é um prazer trabalhar e lidar no dia a dia. Foi um privilégio desenvolver este trabalho 
consigo e tê-lo como mentor. Espero que este não seja o último trabalho que desenvolvamos 
juntos e que possa continuar a aprender consigo no futuro.

Agradeço também ao Professor Jorge Sampaio pela sua disponibilidade ao aceitar ser orientador 
interno para este projeto e pelo tempo dedicado à leitura desta dissertação.

Tenho também de agradecer ao Professor Daniel Galaviz Redondo pelo muito carinho e apoio 
que me deu e tem dado desde que o conheci. Foi sempre e ainda é alguém a quem sei que 
me posso dirigir para aconselhamento e apoio quer académico quer pessoal. Aproveito 
aqui também para mencionar o Professor Manuel Abreu que, juntamente com o Professor 
Daniel, foram coordenadores do mestrado que me encontro agora a acabar. 

Agradeço à Professora Liliana Apolinário pelo seu papel atual como responsável pelo grupo 
Pheno no LIP e também por me ter recebido, orientado e finalmente convencido a seguir os 
meus estudos não só na física das partículas mas especificamente na fenomenologia. Estendo 
este agradecimento aos restantes membros do grupo e ao LIP. 

Tenho aqui de mencionar também a Professora Maria José Varela que, enquanto minha professora 
nos 10.º e 11.º anos, me mostrou que uma continuação de estudos nesta área era o 
que mais se indicava para mim. Agradeço aqui também à Professora Guiomar Evans pelo apoio 
que me deu durante o período de transição de entrada na FCUL e durante toda a licenciatura. 

Passo agora a agradecer ao meu Pai e à minha Mãe por serem quem são e por terem estado e serem figuras 
presentes sempre que precisei. Agradeço também à minha Avó Nana por estar sempre presente com o 
passar dos anos e por ser sempre alguém tão próximo. Agradeço finalmente ao meu Avô António, 
também pela presença assídua e constante. 
Agradeço também às minhas duas tias Susana(s) pelo carinho e aos meus dois primos, 
em ordem de idades, Rodrigo e Bernardo. Relembro aqui também a memória do meu Avô Luís e da 
minha Avó Olga, que foram sempre muito presentes e carinhosos.

À minha Carol agradeço muito por me deixares aturar-te tão assiduamente e continuamente. Aturar-te tanto 
que nem quando te agradeço pela espetacular pessoa que és e sempre para mim foste, consigo ser 
honesto o suficiente para te dizer o quanto te adoro. Também não o posso dizer assim tanto que 
ficas com o ego muito elevado. Gosto muito de ti becas, besi na beh.

Ao careca, açoriano, caríssimo Sr. Ponte mando muitos beijinhos e desejos por muitas, boas e fresquinhas 
jolas: Super Bock de preferência. Temos de combinar voltar ao C5 quando chegarmos aos setenta anos, 
se ele ainda lá estiver.

Vou agora juntar três Miguéis (descobri agora que se escreve com acento, dado ser o plural de uma palavra
acabada em -el). Em primeiro lugar Miguel Monteiro (e com ele David Xia) pelas muitas horas de HOI e de CS 
de muito baixa qualidade; Em segundo Miguel Pinto pelas inúmeras horas daquela infernal cadeira de biomédica 
com aulas às oito; Em terceiro e (desculpa) último lugar, mas não por falta de carinho, Miguel Rodrigues 
que agradeço por mais coisas do que as que cabem nesta página mas que depois de já ter mencionado CS, 
agradeço especialmente pelas muitas horas de grind no MM. Juntamente com ele, e a quebrar a regra das pessoas 
com o mesmo nome, agradeço também muito à Bia pela sua gigante amizade.

De forma a igualar o carinho da sua tese (que tenho aberta neste exato momento), tenho de agradecer muito 
à Bárbara Pereira por ser alguém tão próximo e amigo. És das poucas pessoas que conheci neste curso para quem 
sempre dei \textit{look up to} (espetacular gramática multi-linguística). Agradeço também ao Gabi por nunca 
mais ter feito a barba como daquela vez em que a rapou toda: há pesadelos que espero nunca mais experienciar. 
Tirando isso há pouco a que te agradecer, tirando seres o grande amigo que és e dito de coração, memo um ganda 
tropinha. Juntamente com o Sr. Ponte e o Dr. Pinto, temos nestes agradecimentos reunido todo o gangue do tremoço. 
Vou agradecer aqui também à Vera por estar a manter o Gabi em condições para nós nos últimos meses. 
Para além disso tenho a dizer que também és ganda tropinha. 

Tenho também de agradecer muito, muito, muito à Bia e à Joana, que não precisam aqui de apelidos. Aqui admito 
eu não ter sido a pessoa mais presente, mas vocês foram e ainda são. Acho que ambas sabem que não seria possível 
sem o vosso apoio e a vossa presença. Agradeço também muito ao Hélder que, tal como a Bia e a Joana, conheci muito 
pouco depois de ter entrado nesta faculdade. Foi alguém que vi também crescer ao meu lado e é alguém que tenho 
o muito grande prazer de considerar um amigo.

Não posso também deixar de mencionar o Sr. Tenente David Encarnação, e com ele o resto do gangue do C5 (que evoluiu 
para gangue da sala do núcleo), nomeadamente o Hayden, a Ana, a Raquel, a Filipa, o Edu e o camarada Neves. Agradeço 
também aos meus dois mentorandos, bem como àquele que nunca apareceu. E, já que tal é apropriado, 
menciono outra vez a Raquel que, ainda que não seja minha mentoranda, é como se fosse. 

Finalmente agradeço muito a amizade a todos os que, por puro esquecimento, não foram mencionados. A eles junto também a Faculdade de 
Ciências que em tempos foi jardim do Éden e noutros foi a burocracia do processo de Kafka. Ainda assim, não 
tendo sido a primeira instituição em que entrei, foi aquela em que me mantive, onde me formei e onde conheci a grande 
maioria das pessoas que aqui menciono.
